\chapter{Conclusion}
\label{ch:conclusion}

\section{Summary of Contributions}
\label{sec:conclusion-summary}

This thesis delivered an end-to-end American Sign Language
recognition system that runs entirely on a consumer Android phone
and meets the qualitative and quantitative criteria fixed in the
project proposal. The contributions are summarised below in the
order in which they are developed in the preceding chapters.

\paragraph{Two complementary letter classifiers.}
Chapter~\ref{ch:implementation} described the design and training of
two ASL fingerspelling classifiers that occupy opposite ends of the
size and speed trade-off. The image-based EfficientNet-B0 reaches
\SI{100}{\percent} Top-1 accuracy on the 13\,050-image ASL Alphabet
test partition; the landmark-based MLP, which consumes the
63-dimensional MediaPipe Hands output, reaches \SI{99.38}{\percent}
Top-1 on the detector-recoverable subset of the same partition while
shipping in a \SI{0.24}{MB} checkpoint, \(\sim\)65$\times$ smaller
than the image network. Section~\ref{sec:results-qualitative}
demonstrated that on a budget Android device the relative ordering
of the two classifiers \emph{inverts}: the MLP retains
\SI{100}{\percent} accuracy on both test devices while the image
network drops to \SI{92.9}{\percent} on the entry-level TECNO due to
camera-pipeline sensitivity. The thesis therefore makes a concrete
operational argument for shipping a landmark-based variant on the
letter task, not just on the word task where the literature already
endorses it.

\paragraph{A word classifier that clears the acceptance threshold
with a hands-only feature set.}
Chapter~\ref{ch:implementation} described a five-seed bidirectional
LSTM trained on a two-handed 128-dimensional landmark feature stream,
with sequence augmentation, mixup, and a long-tail weighted sampler.
The five-seed logit-averaging ensemble reached \SI{56.22}{\percent}
Top-1 and \SI{81.09}{\percent} Top-5 on the WLASL-100 test partition
(Table~\ref{tab:results-words-seeds}), clearing the proposal's hard
acceptance threshold of \SI{40}{\percent} Top-1 by \SI{16}{\percent}
percentage points and reaching within four points of the stretch
target of \SI{60}{\percent}. Section~\ref{sec:results-words-literature}
positioned this number in the WLASL-100 literature: a hands-only
landmark model competitive with pose-based methods (Pose-TGCN,
SPOTER) at a fraction of their input footprint.

\paragraph{An incremental ablation that explains the gain from
\SI{20}{\percent} to \SI{60}{\percent}.}
The V2 word classifier was not produced by a single design choice but
by a six-step trajectory documented in
Table~\ref{tab:results-ablation}. The dominant single contribution is
the data-and-architecture bundle that introduced WLASL frame-window
cropping, two-hand extraction with presence flags, and sequence-level
augmentation (+27.33 validation points). The
\texttt{WeightedRandomSampler} (+6.99 points), the increase in mixup
intensity from $\alpha=0.2$ to $\alpha=0.4$ (+4.53 points), and the
five-seed ensemble (+1.23 validation points, +4.08 test points
relative to the per-seed mean) closed the gap to the stretch
target. The ablation is reusable as a recipe for related landmark-based
sequence classification problems on small datasets.

\paragraph{An on-device Android application that exposes all four
classifiers.}
Chapter~\ref{ch:implementation} described a CameraX-driven Android
application that runs the four classifiers through TensorFlow Lite
with the XNNPACK CPU delegate, displays a Top-3 prediction panel and
a self-reported FPS/latency badge, and recovers from MediaPipe's
cold-start timestamp anomaly through a volatile guard around
\texttt{detectAsync}. The Words tab implements the single-shot lock
state machine described in Section~\ref{sec:impl-android-pipeline}:
a \(\sim\)1~s pre-collection countdown, single inference after the
buffer fills, on-screen latching of the result, and automatic
release on an eight-frame hand-down signal. No cloud inference, no
network connectivity, and no server-side component is required at
inference time.

\paragraph{An eight-cell on-device performance characterisation.}
Section~\ref{sec:results-on-device} reported FPS and latency for
each of the four classifiers on two phones: the entry-level
TECNO SPARK 20C (Helio G36) and the flagship Xiaomi 15T Pro
(Dimensity 9400+). The headline result is that, after the
GPU delegate was removed for cold-start stability, the cost of
CPU-only execution is invisible on the flagship (Ensemble at
\SI{30}{FPS} / \SI{80}{ms}) and prohibitive on the entry-level
device (Ensemble at \SI{17}{FPS} / \SI{665}{ms}). The matrix
produces a concrete deployment guideline: on entry-level
hardware only the landmark MLP is interactive; on flagship
hardware all four classifiers are usable and model selection
reduces to the offline accuracy / size trade-off.

\paragraph{A two-device manual evaluation that bounds the offline
{$\rightarrow$} on-device gap.}
Section~\ref{sec:results-qualitative} reported a one-trial-per-class
manual test of the letter classifiers on both devices (29 trials per
cell, 28 measurable) and of the word classifiers on the Xiaomi (100
trials per model). The single-user OOD gap relative to the offline
test partitions was at most seven percentage points for the letter
classifiers and \SIrange{7}{14}{\percent} points for the word
classifiers, with both word classifiers clearing their manual-test
acceptance thresholds (\SI{38.0}{\percent} single, \SI{49.0}{\percent}
ensemble) and a structured set of attractor classes
(\texttt{white}, \texttt{walk}, \texttt{drink}, \texttt{full})
identified as the dominant failure mode. The manual test therefore
provides the empirical evidence behind the deployability claim that
motivates the thesis.

\section{What Was Not Done}
\label{sec:conclusion-not-done}

A clear statement of what falls outside the scope of the present
work is as important as a statement of what was achieved, both for
honest reporting and because the senior-project guidelines require
it. The boundaries listed here are deliberate; each was fixed at
the proposal stage in
Chapter~\ref{ch:introduction} and observed throughout the
implementation and evaluation chapters.

\begin{itemize}
  \item \textbf{Larger WLASL subsets were not evaluated.}
        Word-level training and evaluation are restricted to the
        WLASL-100 subset. The larger 300, 1{,}000, and 2{,}000-class
        subsets of the WLASL benchmark are not addressed; the
        signer-disjoint test protocol, the bidirectional LSTM
        backbone, and the on-device export pipeline are evaluated
        only at the 100-class operating point.
  \item \textbf{Sign languages other than ASL were not considered.}
        Only American Sign Language is recognised. No data, no
        labels, and no model from any other regional sign language
        (e.g.\ T\"urk \.I\c{s}aret Dili, British Sign Language,
        Deutsche Geb\"ardensprache) enters the pipeline. The
        landmark-based front end is in principle language-agnostic,
        but cross-language transfer is not investigated.
  \item \textbf{Continuous sign language translation was not
        attempted.} Both pipelines classify a single isolated
        unit---a static letter or a pre-segmented word
        clip---and produce a class label, not a gloss sequence or a
        translation into written language. Co-articulation between
        adjacent signs, automatic temporal segmentation of a
        free-flowing signing stream, and the mapping from gloss to
        natural-language sentences are out of scope.
  \item \textbf{No cloud or server-side inference component was
        built.} All recognition is performed on the consumer
        Android device, using the exported TensorFlow Lite
        artefacts. The system does not call any remote API at
        inference time and does not require network connectivity.
        Server-side training, ensembling, or re-ranking---common
        in commercial accessibility tools---is therefore not part
        of the deliverable.
  \item \textbf{Signer-level analysis was not performed.}
        The ASL Alphabet dataset does not expose signer identity in
        its public release, and the WLASL-100 evaluation follows
        the dataset's published signer-disjoint split as a single
        fixed partition. Per-signer accuracy breakdowns,
        cross-signer adaptation, and few-shot personalisation to a
        new signer are not reported.
  \item \textbf{User-study evaluation with deaf or
        hard-of-hearing participants was not conducted.} The
        present evaluation is purely technical: classification
        accuracy on held-out data and on-device throughput on the
        target hardware. A field study with members of the
        intended user community---which would be the appropriate
        next step before any deployment claim---was not within the
        scope of a one-semester senior-project thesis and is left
        to future work.
\end{itemize}

\section{Limitations}
\label{sec:conclusion-limitations}

The previous section listed what was placed outside the scope of
the thesis by design. This section lists the residual limitations
of \emph{what was actually built}: properties of the delivered
system that constrain how its measurements should be interpreted
and where its applicability ends. They are grouped into
architectural limitations, which follow from design choices,
data-side limitations, which follow from the public datasets the
system was trained on, and performance-side limitations, which
follow from the on-device measurements of
Section~\ref{sec:results-on-device} and the manual evaluation of
Section~\ref{sec:results-qualitative}.

\paragraph{Hands-only landmark stream.}
Both landmark-based classifiers consume the 21~hand-landmark
output of MediaPipe Hands and nothing else. ASL signs frequently
carry meaning in components that this representation discards:
body location (signs anchored to the chest, shoulder, or face),
upper-arm motion (which a two-handed sign uses to coordinate its
hands), and non-manual markers such as eyebrow position or mouth
shape. The error analysis of
Section~\ref{sec:results-errors} attributes the most frequent
WLASL-100 confusion pairs precisely to this missing context: the
five top pairs all share handshape and gross hand motion and
differ only in body location or facial marker. A model that adds
body pose, and especially facial landmarks, would resolve a
fraction of these errors at near-zero inference cost on the
target hardware (the corresponding extension is the first item of
Section~\ref{sec:conclusion-future}).

\paragraph{Single-hand letter pipeline.}
The landmark-based fingerspelling MLP is trained on the dominant
hand only; if MediaPipe detects two hands in the frame, only the
first is consumed and the rest is discarded. This is appropriate
for the bulk of the ASL Alphabet---most letters use a single
hand---but it leaves regional two-hand variants and the
auxiliary classes under-supported. The image-based EfficientNet
classifier does not have this limitation, since it operates on
raw pixels, but it correspondingly cannot benefit from the
explicit two-handed structure that the LSTM word model uses.

\paragraph{MediaPipe detection floor.}
Both landmark pipelines are gated on MediaPipe's hand detector.
On the ASL Alphabet test split, MediaPipe failed to recover a
hand on \SI{26.91}{\percent} of test images
(Section~\ref{sec:impl-mlp}), including the entire
\texttt{nothing} class and a thumb-tucked subset of \textit{M}
and~\textit{N}. The Android demo handles these frames gracefully
by greying out the landmark-based output, but the practical
implication is the same: under conditions where MediaPipe fails
(low-resolution input, motion blur, partial occlusion, unusual
lighting), the landmark pipelines simply do not produce a
prediction. The image-based pipeline is the appropriate fallback
on such frames; the system does not currently switch between the
two automatically.

\paragraph{ASL Alphabet generalisation cannot be inferred from
the reported test accuracy.}
The \SI{100}{\percent} EfficientNet and \SI{99.38}{\percent} MLP
numbers reported in Section~\ref{sec:results-letters} are
upper-bound figures, obtained on a random 70/15/15 split of a
single-source Kaggle dataset whose images were collected by one
contributor in a controlled indoor setting. The dataset does not
expose signer identity in its public release, so a
signer-disjoint split is not constructible; the random split
allows the same person, lighting, and background to appear in
both partitions. The reported accuracy therefore measures
in-distribution discrimination, not generalisation to a new
signer in a new room with a new camera. The real-device
measurements of Section~\ref{sec:results-qualitative} are the
appropriate ground truth for the latter question.

\paragraph{WLASL-100 person-mixed split.}
The WLASL-100 evaluation uses the dataset's published canonical
split, which is person-mixed: the same signer can appear in both
the train and the test partitions. This is a known property of
the WLASL benchmark and is shared by all published baselines in
Table~\ref{tab:results-wlasl-lit}; the relative comparisons in
that table are therefore on equal terms. The absolute number,
however, is inflated relative to a stricter signer-disjoint
protocol, and a system that reaches \SI{56.22}{\percent} on the
mixed split should not be assumed to reach the same number on
data from a signer that the training set has never seen.

\paragraph{Isolated-only word recognition.}
The word classifier expects a temporally pre-segmented clip
containing exactly one sign. It does not detect when a sign
starts or ends in a continuous signing stream, and it does not
model co-articulation between adjacent signs. The Android demo
applies the same fixed 32-frame window logic at inference time
and therefore inherits this assumption: the user must produce
each word in isolation for the displayed classification to be
meaningful. A real accessibility tool would have to add a
segmentation stage in front of the classifier, an extension that
is identified as a high-impact item in
Section~\ref{sec:conclusion-future}.

\paragraph{Asymmetric latency cost of CPU-only execution on
entry-level hardware.}
To eliminate a cold-start \texttt{MediaPipeException} in the
\texttt{LIVE\_STREAM} path of MediaPipe Hands, the deployed Android
build runs all four classifiers on the CPU through XNNPACK and does
not enable the TensorFlow Lite GPU delegate
(Section~\ref{sec:results-on-device-cpu}). The cost of this decision
is invisible on flagship hardware---on the Xiaomi 15T Pro all four
classifiers ran at \SI{30}{FPS} with end-to-end latencies of at most
\SI{90}{ms}---but it is prohibitive on entry-level hardware. On the
TECNO SPARK 20C, EfficientNet-B0 and the BiLSTM single model both
exceeded \SI{200}{ms} mean latency and the five-seed ensemble
reached \SI{665}{ms} (Table~\ref{tab:on-device-words}), at which
point the application is no longer interactive. The trade-off taken
for build stability is therefore not uniform across the deployment
spectrum: it is free on flagship phones and disqualifying for the
ensemble on budget phones. A future build that exposes a per-device
delegate choice (GPU on devices where the cold-start path is stable,
CPU elsewhere) would partly recover this cost, but the present
thesis does not include that work.

\paragraph{Camera-pipeline sensitivity of the image-based letter
classifier.}
The manual test
(Section~\ref{sec:results-qualitative-letters}) revealed a behavior
that the offline test partition did not. EfficientNet-B0, which
reaches \SI{100}{\percent} on the offline letter test set, drops to
\SI{92.9}{\percent} on the entry-level TECNO and stays at
\SI{100}{\percent} on the Xiaomi. The two errors on the TECNO
(\textit{D}~$\rightarrow$~\textit{K/R} and
\textit{K}~$\rightarrow$~\textit{W}) involve hand shapes that the
landmark MLP classified correctly on the same device with full
confidence; their probable cause is the budget phone's image-signal
processing pipeline (sharper highlights, harsher white-balance
correction, slight motion blur) interacting with a pixel-based
backbone that was trained on the Kaggle ASL Alphabet dataset's
controlled-indoor imagery. The landmark abstraction discards
precisely this kind of camera-conditioned variation. The practical
consequence is that the image classifier's offline-vs-deployed
accuracy gap is itself device-conditioned, and a deployment that
guarantees image-network accuracy across the budget tier would
require either a camera-aware retraining set or a runtime fallback
to the landmark variant.

\paragraph{Single-user OOD ceiling on the manual evaluation.}
The manual evaluation of Section~\ref{sec:results-qualitative} was
conducted by a single human signer (the author). One signer cannot
sample the population of hand morphologies, signing styles,
fingerspelling speeds and ambient camera conditions that a real
deployment would encounter, so the manual-test accuracy numbers are
not representative samples of production accuracy in either
direction. They are useful for two narrower purposes: verifying
that the proposal's qualitative criterion (``the recognised letter
or word appears on screen'') is satisfied in real use, and bounding
the offline-to-on-device gap on a single OOD draw. A
representative production accuracy number would require a
multi-signer evaluation that the present scope did not include,
and is the user-study item listed in
Section~\ref{sec:conclusion-not-done} and revisited as the highest-impact
follow-up in Section~\ref{sec:conclusion-future}.

\section{Future Work}
\label{sec:conclusion-future}

The scope boundaries and limitations identified above suggest a
number of concrete extensions that would be natural follow-ups to
the present thesis. They are listed in increasing order of effort.

\paragraph{Augmenting the landmark front end with body pose.}
The current word-level pipeline consumes only the 21~hand landmarks
of MediaPipe Hands, dropped to zero when the hand leaves the frame.
WLASL signs, however, frequently involve the upper body: signs are
located in space relative to the chest, the dominant hand often
references the non-dominant arm, and two-handed signs depend on
coordinated motion of both hands. Combining MediaPipe Hands with
the body-pose joints of MediaPipe Pose---as Pose-TGCN and the
pose-guided RGB fusion of Hosain et~al.~\cite{hosain2021fusion} do at
the body level---would add this missing context at a near-zero
inference cost on the mobile target, because both extractors
already run on-device. A direct comparison between the present
hands-only landmark stream and a Pose+Hands stream, with the rest
of the architecture held fixed, is the most immediate extension.

\paragraph{Replacing the recurrent temporal model with a
transformer encoder.}
Section~\ref{sec:related-temporal} surveyed the move from recurrent
to transformer-based temporal models on WLASL, with the Sign
Pose-Based Transformer of Boh\'a\v{c}ek and
Hr\'uz~\cite{bohacek2022spoter} reaching 63.18\% Top-1 on the
100-class subset using pose-only inputs. A natural follow-up to the
present bidirectional LSTM is to replace its recurrent backbone
with a comparably sized pose transformer encoder, training and
evaluating it under the same signer-disjoint protocol and the same
on-device export pipeline. The interesting question is not whether
a transformer can match an LSTM on this dataset---the prior
literature suggests it can---but whether the resulting model
remains within the mobile size and latency budget that motivates
this thesis.

\paragraph{Two-handed handling on the letter task.}
The landmark-based fingerspelling MLP is trained on a single hand
and falls back to a zero-vector when no hand is detected. Letters
that involve two hands in some regional variants, and several of
the auxiliary classes, would benefit from a model that consumes
both hands when both are present. This is a small architectural
change---doubling the input dimension and zero-padding the
secondary hand---but it would require re-training on a dataset that
preserves both-hand annotations.

\paragraph{Real-time word boundary detection.}
The word-level pipeline is evaluated under the WLASL-100 protocol,
which assumes pre-segmented isolated clips. A usable accessibility
tool would have to detect the start and end of each sign in a
continuous signing stream. A lightweight on-device boundary
detector---e.g.\ a binary classifier over short windows of the
same landmark stream, or an energy-based segmenter over hand
velocity---could be inserted in front of the existing word
classifier without changing the latter. This extension is mostly
engineering rather than research, but it is the single change with
the largest impact on perceived usability.

\paragraph{Cross-lingual transfer.}
Because the landmark front end is independent of the spoken
language of the signer, the same MLP and LSTM architectures could
in principle be re-trained on a Turkish Sign Language
fingerspelling and isolated-sign dataset. Investigating how much
of the WLASL-100 trained weight can be reused as a warm start,
versus how much must be re-learned from scratch, would address a
practical question for Turkish-speaking deaf communities and would
be a natural continuation of this thesis in a graduate-level
follow-up.

\paragraph{User study with the intended community.}
Finally, all of the above extensions are technical. The single
extension that would have the highest impact on whether the system
is actually useful is a field evaluation with deaf and
hard-of-hearing participants. Such a study would measure not
classification accuracy but task success---can the participant
fingerspell a name and have the application recognise it, can a
hearing interlocutor follow a 100-word vocabulary exchange---and
would produce design feedback (camera framing prompts, confidence
display, error-recovery affordances) that no offline benchmark can
provide.
